\section{Four-factor content and cross-factor pruning}
\label{sec:tpc30-cross-content}

The full target gcd can be read from any pair of target
factorizations, but it is not usually the product of the two
within-side gcds alone.

\begin{proposition}[Exact four-factor content decomposition]
\label{prop:tpc30-four-factor-content}
Let
\begin{equation}
 N_m=rU,\qquad N_n=sV
 \label{eq:tpc30-two-factorizations}
\end{equation}
with positive integers $r,s,U,V$.  Put
\begin{equation}
 g=(r,s),\qquad e=(U,V),
 \label{eq:tpc30-within-contents}
\end{equation}
and write
\begin{equation}
 r=gr_0,\quad s=gs_0,\qquad
 U=eU_0,\quad V=eV_0.
 \label{eq:tpc30-primitive-factorizations}
\end{equation}
Thus $(r_0,s_0)=(U_0,V_0)=1$.  Define the two normalized
cross-contents
\begin{equation}
 x=(r_0,V_0),\qquad y=(s_0,U_0).
 \label{eq:tpc30-cross-contents}
\end{equation}
Then
\begin{equation}
 \boxed{
 (N_m,N_n)=gexy.}
 \label{eq:tpc30-four-factor-identity}
\end{equation}
In particular, on the primitive physical orbit,
\begin{equation}
 \boxed{
 gexy=\cG_{m,n}(j)\mid m-n.}
 \label{eq:tpc30-four-factor-determinant}
\end{equation}
No coprimality between $g$ and $e$ is asserted or required.
\end{proposition}

\begin{proof}
Factoring the two common within-side contents gives
\begin{equation}
 (N_m,N_n)
 =ge\,(r_0U_0,s_0V_0).
 \label{eq:tpc30-remove-ge}
\end{equation}
We prove
\begin{equation}
 (r_0U_0,s_0V_0)
 =(r_0,V_0)(s_0,U_0)
 \label{eq:tpc30-cross-gcd-identity}
\end{equation}
prime by prime.  Fix a prime $p$ and write
\[
 a=v_p(r_0),\quad b=v_p(s_0),\quad
 u=v_p(U_0),\quad v=v_p(V_0).
\]
The two primitive conditions give
\[
 \min\{a,b\}=0,\qquad \min\{u,v\}=0.
\]
If $a>0$, then $b=0$.  When $u>0$, one has $v=0$ and both sides of
\begin{equation}
 \min\{a+u,b+v\}
 =\min\{a,v\}+\min\{b,u\}
 \label{eq:tpc30-valuation-identity}
\end{equation}
are zero; when $v>0$, one has $u=0$ and both sides are
$\min\{a,v\}$.  If $u=v=0$, both sides are again zero.  The case
$b>0$ is symmetric, and the case
$a=b=0$ is immediate.  This proves
\eqref{eq:tpc30-cross-gcd-identity}.  Equations
\eqref{eq:tpc30-remove-ge} and
\eqref{eq:tpc30-target-gcd-identity} now give the result.
\end{proof}

\begin{remark}[Why the cross terms cannot be dropped]
\label{rem:tpc30-cross-terms-needed}
Take $r=2$, $U=3$, $s=3$, and $V=4$.  Then $g=e=1$, but
\[
 (rU,sV)=(6,12)=6.
\]
The missing factors are exactly
$x=(2,4)=2$ and $y=(3,3)=3$.  Thus the generally valid statement is
not $(N_m,N_n)=ge$, but the four-factor identity
\eqref{eq:tpc30-four-factor-identity}.
\end{remark}

\subsection{Expanded divisor channels}

Each kernel in \eqref{eq:tpc30-three-raw-kernels} is a finite sum over
one divisor of each target.  In a mixed channel, one selected divisor
comes from the prefix and the other from the ultra complement; in the
both-ultra channel both are reflected complement variables.  For any
chosen pair $r\mid N_m$, $s\mid N_n$, put
\begin{equation}
 U=N_m/r,\qquad V=N_n/s.
 \label{eq:tpc30-co-divisors}
\end{equation}
Indeed, for a prefix leg and a reflected ultra leg, respectively,
\begin{align}
 \sum_{r\mid N_m}|b_R(r)|
 &\ll \tau(N_m)(\log(2X))^3,
 \label{eq:tpc30-prefix-absolute-mass}\\
 \sum_{r\mid N_m}|a(N_m/r)|
 &\ll \tau(N_m)\log(2X).
 \label{eq:tpc30-ultra-absolute-mass}
\end{align}
The same estimates hold for $N_n$.  Since both targets are
$O_{\mathscr D}(X)$, the divisor bound shows that the absolute product
mass of any mixed or both-ultra subfamily is bounded by the full
pointwise divisor envelope.  Here $\kappa^{(\nu)}(j,r,s)$ denotes the
product of the two divisor coefficients in the expanded $\nu$-th
kernel, so
\begin{equation}
 \sum_{\substack{r\mid N_m,\ s\mid N_n\\
                  \textup{retained channel supports}}}
 |\kappa^{(\nu)}(j,r,s)|
 \ll_{\eps,\mathscr D}X^\eps.
 \label{eq:tpc30-expanded-envelope}
\end{equation}
Consequently the proof of
\cref{thm:tpc30-common-target-pruning} applies unchanged after
restricting the expanded divisor sum to any condition implying
$\cG_{m,n}(j)>Z$.

\begin{corollary}[Cross-factor pruning]
\label{cor:tpc30-cross-factor-pruning}
For each mixed or both-ultra expanded channel, restrict its individual
divisor terms by any one of
\begin{equation}
 ge>Z,\qquad e>Z,\qquad x>Z,\qquad y>Z,
 \qquad\text{or more generally}\qquad gexy>Z.
 \label{eq:tpc30-large-cross-events}
\end{equation}
The resulting physical contribution satisfies
\begin{equation}
 \boxed{
 |\cS_{\rm cross>Z}|
 \ll_{\eps,h,\mathscr D}X^\eps
 \left(Q^2+\frac{XQ}{Z}\right).}
 \label{eq:tpc30-cross-factor-bound}
\end{equation}
If $Z\ge X^\eta$, every fixed
$\eta_0<\min\{\eta,1-\beta\}$ is a power-saving exponent.
\end{corollary}

\begin{proof}
By \cref{prop:tpc30-four-factor-content}, each of the first four
conditions in \eqref{eq:tpc30-large-cross-events} implies
$\cG_{m,n}(j)>Z$, while the last is equivalent to that condition.
Use \eqref{eq:tpc30-expanded-envelope} in place of
\eqref{eq:tpc30-kernel-envelope} in the proof of
\cref{thm:tpc30-common-target-pruning}.
\end{proof}

\subsection{What is new relative to reflected content}

TPC-29 tracks
\[
 g=(r,s),\qquad q=[r,s],
\]
and proves a power-saving bound for a long-fiber sector with large $g$
and a sparse sector in
which the primitive reflected lcm $q/g$ is small.  The factors
\[
 e=(U,V),\qquad
 x=(r/g,V/e),\qquad
 y=(s/g,U/e)
\]
are invisible if the two selected divisor variables are considered
without their co-divisors.  Thus
\cref{cor:tpc30-cross-factor-pruning} can be effective with $g=1$ and
$q/g$ large, exactly where the preceding absolute reflected-incidence
argument reaches only the natural scale.

No squarefreeness of $r$ or $s$ is used.  In an ultra term the
co-divisor coefficient $a(U)$ or $a(V)$ vanishes unless the relevant
co-divisor is squarefree, but the four-factor identity itself is an
integer identity and requires no coefficient support.
